\begin{thebibliography}{99}
\bibitem{stinespring}
W. F. Stinespring, Positive functions on $C^*$-algebras,
\emph{Proc. Amer. Math. Soc.} \textbf{6}, 211--216 (1955).
\bibitem{connes}
A. Connes, Classification of injective factors. Cases $\mathrm{II}_1$, $\mathrm{II}_\infty$,
$\mathrm{III}_\lambda$, $\lambda\ne1$,
\emph{Ann. Math.} \textbf{104}, 73--115 (1976).
\bibitem{hm11}
U. Haagerup and M. Musat, Factorization and dilation problems for completely positive maps
on von Neumann algebras, \emph{Commun. Math. Phys.} \textbf{303}, 555--594 (2011),
\href{https://arxiv.org/abs/1009.0778}{arXiv:1009.0778}.
\bibitem{hm15}
U. Haagerup and M. Musat, An asymptotic property of factorizable completely positive maps
and the Connes embedding problem, \emph{Commun. Math. Phys.} \textbf{338}, 721--752 (2015),
\href{https://arxiv.org/abs/1408.6476}{arXiv:1408.6476}.
\bibitem{mr21}
M. Musat and M. R\o rdam, Factorizable maps and traces on the universal free product of
matrix algebras, \emph{Int. Math. Res. Not. IMRN} \textbf{2021}(23), 17951--17970 (2021),
\href{https://arxiv.org/abs/1903.10182}{arXiv:1903.10182}.
\bibitem{musat}
M. Musat, The Connes--Kirchberg problem and infinite-dimensional phenomena in quantum
information theory, \href{https://arxiv.org/abs/2512.00432v1}{arXiv:2512.00432v1} (2025), to appear in
\emph{M\"unster J. Math.}
\bibitem{mip}
Z. Ji, A. Natarajan, T. Vidick, J. Wright, and H. Yuen,
$\mathrm{MIP}^*=\mathrm{RE}$, \href{https://arxiv.org/abs/2001.04383v3}{arXiv:2001.04383v3} (2022).
\bibitem{openai}
OpenAI, Nonsofic groups exist, Chapter~3 of \emph{Ten Advances in Mathematics and Theoretical
Computer Science} (2026), \url{https://cdn.openai.com/pdf/ten-proofs-oai.pdf}.
\bibitem{kt}
G. Kun and A. Thom, Nonsofic wreath products of residually finite groups,
\href{https://arxiv.org/abs/2608.06222v3}{arXiv:2608.06222v3} (2026), Theorems A, E and 4.1.
\bibitem{at}
V. Alekseev and A. Thom, Centralizers of sofic approximations of Kazhdan groups,
\href{https://arxiv.org/abs/2608.05362}{arXiv:2608.05362} (2026), Open Problem 6.2(a).
\bibitem{thom}
A. Thom, A conditional construction of a nonhyperlinear group and the centralizer problem,
\href{https://arxiv.org/abs/2609.27795v1}{arXiv:2609.27795v1} (14 September 2026), preprint,
Theorems 1.2 and 1.3, Remark 4.4, and Section 5.2.
\bibitem{liu}
J. Liu, Nonhyperlinear groups exist, manuscript dated 21 September 2026 (first version posted
20 September 2026), formerly at
\nolinkurl{https://jihaoliu.org/ai-results/nonhyperlinear-groups-exist-2026-09-20.pdf}, archived on
26 September 2026 at
\url{https://web.archive.org/web/20260926151036/https://jihaoliu.org/ai-results/nonhyperlinear-groups-exist-2026-09-20.pdf},
Theorems 1.1 and 1.2.
\bibitem{alt}
V. Alekseev, J. Liu, and A. Thom, Spectral gap, internality of relative commutants, and
non-hyperlinear groups, \href{https://arxiv.org/abs/2609.37761v1}{arXiv:2609.37761v1}
(29 September 2026), preprint, Theorem B, Corollary C, and Proposition 2.2.
\bibitem{higman}
G. Higman, Subgroups of finitely presented groups, \emph{Proc. Roy. Soc. London Ser. A}
\textbf{262}, 455--475 (1961).
\bibitem{ejz}
M. Ershov and A. Jaikin-Zapirain, Property $(T)$ for noncommutative universal lattices,
\emph{Invent. Math.} \textbf{179}, 303--347 (2010),
\href{https://arxiv.org/abs/0809.4095v2}{arXiv:0809.4095v2}, Theorems 1.2 and 6.2 and Corollary 4.7.
\bibitem{mr}
M. Musat and M. R\o rdam, Non-closure of quantum correlation matrices and factorizable
channels that require infinite dimensional ancilla (with an appendix by N. Ozawa),
\emph{Commun. Math. Phys.} \textbf{375},
1761--1776 (2020), \href{https://arxiv.org/abs/1806.10242}{arXiv:1806.10242}.
\bibitem{mori}
M. Mori, On the shape of correlation matrices for unitaries, \emph{Math. Scand.}
\textbf{130}, 359--363 (2024), \href{https://doi.org/10.7146/math.scand.a-142800}{doi:10.7146/math.scand.a-142800},
\href{https://arxiv.org/abs/2308.03345}{arXiv:2308.03345}.
\bibitem{ad}
C. Anantharaman-Delaroche, On ergodic theorems for free group actions on noncommutative spaces,
\emph{Probab. Theory Related Fields} \textbf{135}, 520--546 (2006).
\bibitem{pr}
V. I. Paulsen and M. Rahaman, Bisynchronous games and factorizable maps,
\emph{Ann. Henri Poincar\'e} \textbf{22}, 593--614 (2021),
\href{https://arxiv.org/abs/1908.03842}{arXiv:1908.03842}.
\bibitem{bhv}
B. Bekka, P. de la Harpe, and A. Valette, \emph{Kazhdan's Property (T)}, New Mathematical
Monographs \textbf{11} (Cambridge University Press, 2008).
\bibitem{slofstra}
W. Slofstra, Tsirelson's problem and an embedding theorem for groups arising from non-local
games, \emph{J. Amer. Math. Soc.} \textbf{33}, 1--56 (2020),
\href{https://arxiv.org/abs/1606.03140}{arXiv:1606.03140}.
\bibitem{ozawa}
N. Ozawa, About the QWEP conjecture, \emph{Internat. J. Math.} \textbf{15}, 501--530 (2004),
\href{https://arxiv.org/abs/math/0306067}{arXiv:math/0306067}.
\bibitem{causal}
S. Douglas and N. Mghirbi, Causal quantum-channel simulation: memory beyond entropy,
version 1.1.2 (2026), \href{https://doi.org/10.5281/zenodo.23070870}{doi:10.5281/zenodo.23070870},
Definition 2.3, Theorem 3.2, Lemmas 5.4 and 5.9, and Theorem 5.8.
\bibitem{configuration}
S. Douglas and N. Mghirbi, Amenable groups with nearly exponential sofic profile, and quantum
channels that need nearly linear memory, version 1.0.2 (2026),
\href{https://doi.org/10.5281/zenodo.23070871}{doi:10.5281/zenodo.23070871},
Proposition 8.7 and Lemma 8.8.
\end{thebibliography}
